\documentclass[conference]{IEEEtran}
\usepackage[T1]{fontenc}
\usepackage[utf8]{inputenc}
\usepackage[brazil]{babel}
\usepackage{amsmath}
\usepackage{amssymb}
\usepackage{booktabs}
\usepackage{array}
\usepackage{graphicx}
\usepackage{url}
\usepackage{hyperref}
\usepackage{cite}

\hypersetup{
    colorlinks=true,
    linkcolor=blue,
    urlcolor=blue
}

\begin{document}

\title{Tarefa Avaliativa 1 --- Sistemas de Comunicação 2026/1}

\author{
  \IEEEauthorblockN{Bruno Becker Silva}
  \IEEEauthorblockA{Escola Politécnica --- PUCRS\\
  Disciplina: Sistemas de Comunicação\\
  Professor: Flavio Eduardo Soares e Silva}
}

\maketitle

% ============================================================
\section*{Questão 1 — O Sistema de Comunicação como Cadeia Integrada}
% ============================================================

A comunicação de dados não pode ser tratada como um fenômeno isolado em um único componente,
pois envolve uma cadeia causal em que cada elemento condiciona o comportamento dos demais.
O modelo clássico, derivado de Shannon e Weaver (1948) e consolidado na literatura de Rochol~\cite{rochol2012},
organiza essa cadeia em cinco blocos funcionais: \textbf{fonte}, \textbf{transmissor}, \textbf{meio},
\textbf{receptor} e \textbf{destino}. Cada bloco impõe restrições distintas sobre desempenho,
confiabilidade e custo do sistema como um todo.

\subsection*{Papel de cada elemento e suas limitações}

A \textbf{fonte} é responsável por gerar a informação e, portanto, determina a taxa de entropia
do sinal: quanto maior a variabilidade da informação, maior a largura de banda mínima exigida
do canal. Uma fonte mal projetada --- por exemplo, sem compressão adequada --- desperdiça
capacidade do canal e eleva o custo operacional.

O \textbf{transmissor} adapta o sinal da fonte ao meio físico, realizando operações de
codificação de linha, modulação e, quando necessário, codificação de canal para proteção contra
erros. Limitações aqui se traduzem em potência insuficiente, mismatch de impedância ou
modulações inadequadas ao perfil de ruído do meio, degradando a relação sinal-ruído (SNR).

O \textbf{meio de transmissão} é o elemento de maior variabilidade: sua largura de banda,
atenuação por quilômetro, susceptibilidade a interferências eletromagnéticas (EMI) e retardo
de propagação determinam o teto físico da comunicação. Uma limitação no meio não pode ser
compensada apenas por melhorias no transmissor ou receptor, pois o teorema de Shannon estabelece
o limite de capacidade como $C = B \log_2(1 + \text{SNR})$, onde $B$ é a largura de banda
disponível no meio.

O \textbf{receptor} deve recuperar o sinal original a partir de uma versão degradada. A qualidade
de seus circuitos de filtragem, equalização e detecção de limiar define a taxa de erro de bit
(BER) na presença de ruído. Um receptor com alta figura de ruído (\textit{noise figure}) eleva
o piso de ruído efetivo, reduzindo a margem de enlace.

O \textbf{destino} impõe requisitos de latência, jitter e taxa mínima de dados que retroalimentam
o dimensionamento de todos os outros blocos. Em sistemas de vídeo conferência, por exemplo,
latência acima de 150~ms já prejudica a inteligibilidade da conversa, exigindo QoS nos
equipamentos de rede.

\subsection*{Exemplos práticos contemporâneos}

\textbf{Exemplo 1 --- Rede 5G:} Na infraestrutura 5G, a fonte (smartphone) gera fluxos de
vídeo 4K. O transmissor aplica modulação OFDM com até 256-QAM. O meio (ar, frequências entre
700~MHz e 26~GHz dependendo da faixa) impõe atenuação crescente com a frequência e com
obstáculos físicos. O receptor (eNB/gNB) utiliza MIMO massivo para diversidade espacial. Uma
limitação no meio --- como congestionamento em 3,5~GHz por interferência de radar --- força
\textit{handover} para faixas de menor capacidade, reduzindo a taxa efetiva percebida pelo
destino. O custo aumenta proporcionalmente à quantidade de pequenas células necessárias para
compensar as limitações do meio em alta frequência.

\textbf{Exemplo 2 --- Link de fibra óptica intercontinental (submarine cable):} A fonte
(data center) gera tráfego IP agregado. O transmissor aplica WDM (\textit{Wavelength Division
Multiplexing}) com dezenas de comprimentos de onda simultâneos. O meio (fibra monomodo
submersa) apresenta atenuação de $\approx 0{,}2$~dB/km, exigindo amplificadores EDFA a cada
$\sim$80~km. Uma limitação no receptor --- por exemplo, fotodetectores com baixa
responsividade --- aumenta a BER, demandando maior margem de potência e, portanto, mais
repetidores, elevando o custo de implantação do cabo. Uma falha no meio (rompimento mecânico
por âncora) pode interromper a comunicação de países inteiros, evidenciando a criticidade de
cada elo da cadeia.

\subsection*{Síntese}

A análise sistêmica é necessária porque a otimização de um único elemento sem considerar os
demais pode ser ineficaz ou até contraproducente. Um transmissor de altíssima potência não
resolve limitação de banda do meio; um receptor ultrassensível não compensa modulação inadequada
do transmissor. O desempenho, a confiabilidade e o custo emergem da interação entre todos os
cinco elementos.

% ============================================================
\section*{Questão 2 — O Modelo OSI: Camadas, Serviços e Dependências}
% ============================================================

O modelo OSI (\textit{Open Systems Interconnection}) organiza as funções de comunicação em
sete camadas hierárquicas, cada uma resolvendo um problema específico e oferecendo serviços
bem definidos à camada superior, enquanto consome serviços das camadas inferiores. A seguir
analisa-se três camadas: \textbf{Física} (camada 1), \textbf{Enlace} (camada 2) e
\textbf{Rede} (camada 3).

\subsection*{Camada Física (Layer 1)}

\textbf{Problema que resolve:} A camada física trata da representação dos bits como sinais
elétricos, ópticos ou de rádio sobre o meio físico. Define tensões, impedâncias, codificações
de linha (NRZ, Manchester, 4B5B), sincronismo de bit e características mecânicas dos conectores.

\textbf{Serviço à camada superior:} Oferece à camada de Enlace uma interface de transmissão
de bits sem interpretação semântica --- um serviço de \textit{bit pipe}, com transparência
de dados e garantia apenas de que os bits são transmitidos (mas não necessariamente sem erro).

\textbf{Dependência das camadas inferiores:} A camada Física \textit{é} a camada mais baixa do
modelo OSI; sua dependência é diretamente com o meio físico (par trançado, fibra, ar), que não
faz parte do modelo formal mas condiciona todos os parâmetros elétricos e ópticos definidos
nesta camada.

\subsection*{Camada de Enlace (Layer 2)}

\textbf{Problema que resolve:} O Enlace transforma o serviço de bit pipe --- sujeito a erros ---
em um canal lógico confiável entre nós diretamente conectados. Realiza enquadramento
(\textit{framing}), detecção e correção de erros (CRC, ARQ), controle de fluxo e endereçamento
físico (MAC).

\textbf{Serviço à camada superior:} Entrega à camada de Rede quadros íntegros entre nós
adjacentes, abstraindo as particularidades do meio físico. A Rede não precisa conhecer se o
enlace é Ethernet, Wi-Fi ou PPP.

\textbf{Dependência das camadas inferiores:} Depende integralmente da camada Física para a
temporização de bits e da integridade elétrica do meio. Se a Física apresentar alta taxa de
erros (SNR baixo), o mecanismo ARQ do Enlace será sobrecarregado com retransmissões, degradando
o throughput efetivo.

\subsection*{Camada de Rede (Layer 3)}

\textbf{Problema que resolve:} A Rede resolve o problema do roteamento fim-a-fim entre hosts
em redes heterogêneas e geograficamente distribuídas. Realiza endereçamento lógico (IP),
roteamento, fragmentação de pacotes e controle de congestionamento.

\textbf{Serviço à camada superior:} Oferece à camada de Transporte um serviço de entrega de
datagramas entre quaisquer dois hosts na internet, abstraindo a topologia intermediária.

\textbf{Dependência das camadas inferiores:} Depende do Enlace para a entrega hop-a-hop.
Incompatibilidades no MTU (\textit{Maximum Transmission Unit}) do Enlace obrigam a camada
de Rede a fragmentar pacotes, aumentando overhead e latência.

\subsection*{Violação de responsabilidades entre camadas}

Quando uma camada assume funções de outra, ocorrem acoplamentos indevidos que comprometem
interoperabilidade e manutenção. Um exemplo clássico é o protocolo
\textit{TCP over wireless}: originalmente, o TCP interpretava perda de pacotes como sinal
de congestionamento (lógica da camada de Transporte), acionando redução de janela. Em enlaces
sem fio, a perda é predominantemente por erro de bit (responsabilidade do Enlace). O TCP não
foi projetado para distinguir esses casos, pois pressupunha que o Enlace entregaria quadros
sem erro. O resultado é a queda desnecessária de throughput em redes Wi-Fi de alta taxa de
erro, só corrigida por mecanismos como \textit{TCP-Westwood} ou ARQ seletivo na camada 2.
Esse exemplo ilustra como a violação da abstração de camada --- o Transporte precisando
conhecer a natureza do meio --- degrada tanto o desempenho quanto a capacidade de evolução
independente de cada camada.

% ============================================================
\section*{Questão 3 — Análise de Sinais no Tempo e na Frequência}
% ============================================================

\subsection*{Complementaridade das representações temporal e espectral}

As representações no domínio do tempo e no domínio da frequência são relacionadas biunivocamente
pela Transformada de Fourier: $X(f) = \int_{-\infty}^{\infty} x(t)e^{-j2\pi ft}dt$. Isso
significa que \textit{toda informação presente em uma representação está implícita na outra};
não há perda de informação ao transitar entre os dois domínios. A visão temporal mostra como
a amplitude do sinal evolui no tempo --- útil para sincronismo, detecção de transições e
análise de jitter. A visão espectral revela a distribuição de energia pelas frequências ---
essencial para dimensionamento de banda, análise de interferência e projeto de filtros.
As duas visões são, portanto, complementares: cada uma torna explícito o que a outra obscurece.

\subsection*{Efeitos sobre a forma de onda e a recuperação da informação}

\textbf{Aumento da taxa de bits:} Maior taxa implica pulsos mais curtos no tempo. Pelo par
de Fourier (relação de Heisenberg-Gabor), pulsos mais estreitos possuem espectros mais largos.
Assim, dobrar a taxa de bits aproximadamente dobra a largura de banda necessária. Se o canal
não suportar essa banda, os pulsos sofrem alargamento temporal (\textit{ISI --- Inter-Symbol
Interference}), tornando difícil a decisão no receptor sem equalização.

\textbf{Limitação de banda do canal:} Um canal passa-baixa ideal com largura de banda $B$
(Hz) é equivalente, pelo critério de Nyquist, a uma capacidade máxima de $2B$ símbolos por
segundo sem ISI. Se a taxa de símbolos exceder $2B$, os pulsos no domínio do tempo se
alargam e se sobrepõem, causando ISI inevitável. Na forma de onda, isso se manifesta como
transições suavizadas e olhos fechados no diagrama de olho, reduzindo drasticamente a margem
de decisão do receptor.

\textbf{Ruído aditivo (AWGN):} O ruído gaussiano branco adiciona componentes aleatórios em
todas as frequências, reduzindo a SNR. No domínio do tempo, traduz-se em flutuações
estocásticas na amplitude do sinal. A recuperação é estatística: o receptor define um limiar
de decisão e a probabilidade de erro é função de $Q(\sqrt{2E_b/N_0})$, onde $E_b$ é a
energia por bit e $N_0$ é a densidade espectral de ruído. O aumento de energia por símbolo
ou a redução da taxa (para aumentar $E_b$) são estratégias diretas de mitigação.

\textbf{Distorção:} Diferentemente do ruído, a distorção é determinística: o canal modifica
amplitude e fase de forma dependente da frequência (resposta não plana). No domínio do tempo,
provoca ISI mesmo na ausência de ruído, pois diferentes componentes espectrais do pulso
chegam com atrasos distintos (\textit{dispersão de grupo}). A recuperação exige equalização
adaptativa (filtros ZF ou MMSE), que inverte a resposta do canal. A distorção não pode ser
mitigada apenas aumentando a potência de transmissão, ao contrário do ruído.

% ============================================================
\section*{Questão 4 — Par Trançado versus Fibra Óptica}
% ============================================================

A escolha do meio físico é uma decisão de engenharia multidimensional. A
Tabela~\ref{tab:comparacao} sumariza as principais dimensões de comparação.

\begin{table}[h!]
\renewcommand{\arraystretch}{1.3}
\caption{Comparação entre par trançado e fibra óptica}
\label{tab:comparacao}
\centering
\begin{tabular}{>{\bfseries}p{2.2cm} p{2.5cm} p{2.5cm}}
\toprule
\textbf{Critério} & \textbf{Par Trançado (UTP/STP)} & \textbf{Fibra Óptica} \\
\midrule
Atenuação & $\sim$20--30~dB/km (Cat6A a 500~MHz) & $\sim$0,2~dB/km (SM 1310~nm) \\
\midrule
Imunidade a EMI & Moderada (STP) / Baixa (UTP) & Total (dielétrico) \\
\midrule
Banda disponível & Até $\sim$2~GHz (Cat8) & Dezenas de THz (WDM) \\
\midrule
Custo de material & Baixo & Médio a alto \\
\midrule
Custo de instalação & Baixo (conectores simples) & Alto (fusão, testes OTDR) \\
\midrule
Facilidade de manutenção & Alta (conectores RJ-45 padrão) & Requer equipamentos especializados \\
\midrule
Escalabilidade & Limitada por distância e frequência & Praticamente ilimitada (WDM) \\
\midrule
Risco de interceptação & Emite campo EM & Zero (não irradia) \\
\bottomrule
\end{tabular}
\end{table}

\subsection*{Adequação por cenário}

O par trançado (Cat6A ou Cat8) é a escolha racional para cabeamento horizontal em edifícios
corporativos, onde as distâncias são inferiores a 100~m, o custo de instalação deve ser
minimizado e a mão de obra qualificada para fibra é escassa. A substituição de toda a
infraestrutura de acesso por fibra nesses ambientes elevaria o custo de implantação em uma
ordem de grandeza sem benefício proporcional de desempenho para aplicações típicas (1~Gbps).

A fibra óptica monomodo é indispensável em enlaces de \textit{backbone}, data centers e
conexões entre edifícios separados por distâncias superiores a 100~m ou em ambientes com
forte EMI (ambientes industriais, próximos a motores de alta potência). Sua imunidade
eletromagnética elimina qualquer acoplamento indutivo ou capacitivo, garantindo integridade
do sinal independentemente do ambiente elétrico.

Em ambientes industriais com inversores de frequência (que geram ruído de $\sim$kHz a
$\sim$MHz), o par trançado --- mesmo blindado --- pode apresentar degradação de SNR; a fibra
resolve estruturalmente esse problema sem qualquer configuração adicional.

\subsection*{Conclusão}

``Melhor meio'' é uma função do contexto de engenharia: distância, orçamento, ambiente
eletromagnético, taxa de dados requerida e perfil de manutenção da organização. O engenheiro
deve dimensionar o meio pelo problema concreto, não pela capacidade nominal máxima.

% ============================================================
\section*{Questão 5 — Largura de Banda, Taxa de Transmissão e Qualidade do Canal}
% ============================================================

\subsection*{Relações fundamentais}

A taxa máxima de transmissão sem erro em um canal ideal (sem ruído, banda $B$) é dada pelo
critério de Nyquist: $R_{max} = 2B\log_2 M$, onde $M$ é o número de níveis de sinalização.
Na presença de ruído gaussiano branco, Shannon estabelece o limite absoluto:
$C = B\log_2(1 + \text{SNR})$. Nenhuma técnica de codificação ou modulação pode superar
$C$ com probabilidade de erro arbitrariamente pequena.

\subsection*{Limitações físicas do meio versus limitações de projeto}

\textbf{Limitações físicas do meio} são intrínsecas ao material: a atenuação por resistência
ôhmica em condutores metálicos, a dispersão cromática em fibras ópticas, o desvanecimento
multipercurso em canais de rádio e o ruído térmico ($kTB$, onde $k$ é a constante de
Boltzmann). Essas limitações não podem ser eliminadas por projeto --- apenas contornadas
(por exemplo, utilizando amplificação, equalização ou diversidade de antenas).

\textbf{Limitações de projeto} são consequências de escolhas de engenharia: modulação com
eficiência espectral subótima (BPSK quando 64-QAM é viável), ausência de codificação de
canal (FEC), sincronismo mal projetado ou topologia de rede inadequada. Essas limitações
\textit{podem} ser eliminadas por redesign do sistema.

\subsection*{Por que aumentar potência nem sempre resolve}

Aumentar a potência de transmissão eleva a SNR e, portanto, pode reduzir a BER em canais
limitados por ruído AWGN. Contudo, essa estratégia falha em quatro situações:

\begin{enumerate}
\item \textbf{Interferência dominante:} Em sistemas como Wi-Fi, o ruído é frequentemente
dominado por interferência de outros transmissores co-canal (não é gaussiano branco). Mais
potência no transmissor local também aumenta a interferência que ele causa nos vizinhos,
resultando em aumento proporcional do patamar de interferência. A SNR efetiva não melhora
e o sistema pode entrar em colapso por poluição eletromagnética.

\item \textbf{Distorção não linear:} Amplificadores de potência operam de forma linear apenas
abaixo de um ponto de compressão de 1~dB (P1dB). Acima desse ponto, introduzem distorção
harmônica e intermodulação, degradando a constelação de modulação e aumentando a BER, mesmo
com potência elevada.

\item \textbf{Limitação de banda:} Se o problema é ISI por banda insuficiente (distorção
determinística), a potência adicional não resolve, pois o sinal degradado é o resultado da
resposta em frequência do canal, não da SNR. A solução correta é a equalização.

\item \textbf{Regulamentação espectral:} Limites de potência isotrópica radiada efetiva
(EIRP) são impostos por órgãos reguladores (ANATEL, FCC), tornando essa estratégia
inviável em sistemas de rádio sem licença adicional.
\end{enumerate}

% ============================================================
\section*{Questão 6 — Projeto de Enlace Inter-Predial: Dados, Voz e Vídeo}
% ============================================================

\subsection*{Contexto e requisitos}

Interligar dois prédios de um campus com tráfego convergente (dados, voz VoIP e vídeo
conferência) exige análise multidimensional: a voz tolera latência até $\sim$150~ms e perda
de pacotes até 1\%; o vídeo exige jitter controlado ($<$30~ms) e banda mínima garantida;
os dados toleram latência maior, mas exigem confiabilidade na entrega.

\subsection*{Meio físico principal: fibra óptica monomodo}

Para um enlace inter-predial em campus, a \textbf{fibra óptica monomodo} (OS2, 1310/1550~nm)
é a escolha tecnicamente superior por três razões principais: (1) a distância entre prédios
tipicamente excede o limite de 100~m do par trançado; (2) o ambiente externo expõe o cabo a
EMI de fontes como motores, geradores e sistemas de climatização; (3) a capacidade de expansão
via WDM permite multiplicar a capacidade sem substituição do cabo físico.

A solução completa inclui: dois switches de borda com uplinks de fibra (SFP+ 10GbE),
patch panels de fibra LC/PC nos distribuidores de cada prédio, e roteador de borda com
suporte a QoS diferenciado (DSCP).

\subsection*{Requisitos mínimos de integridade do sinal}

O enlace deve apresentar perda óptica total inferior a 3~dB (considerando emendas, conectores
e margem de sistema), SNR suficiente para operação em 10~Gbps, e latência de propagação
inferior a 0,5~ms (para distâncias típicas de campus de até 100~m).

\subsection*{Papel das camadas no tratamento da comunicação}

A Camada Física garante a transmissão de bits com integridade óptica. A Camada de Enlace
(Ethernet 802.3) provê enquadramento e detecção de erros (FCS/CRC-32), além de controle de
fluxo (802.3x). A Camada de Rede (IP) realiza o roteamento entre as sub-redes dos dois
prédios e aplica marcação DSCP para QoS. A Camada de Transporte (UDP para voz/vídeo, TCP
para dados) lida com controle de fluxo fim-a-fim e confiabilidade onde necessário.

\subsection*{Principais riscos técnicos}

Os riscos incluem: (1) \textbf{corte físico do cabo} por obras ou vandalismo --- mitigado
por dutos subterrâneos protegidos e, idealmente, por redundância de rota; (2)
\textbf{degradação de conectores} por contaminação de partículas no campo óptico ---
prevenida por inspeção com microscópio de vídeo antes da instalação; (3)
\textbf{jitter de rede} em tráfego de voz/vídeo por congestionamento --- mitigado por
filas com prioridade (WFQ/CBWFQ) e policiamento de tráfego.

\subsection*{Critérios para expansão futura}

A arquitetura deve prever: portas SFP+ livres nos switches (capacidade de 40/100GbE);
dutos com 40\% de espaço livre para cabos adicionais; e documentação OTDR de referência
para comparação futura. A adoção de WDM passivo pode quadruplicar a capacidade sem novo
cabo.

% ============================================================
\section*{Questão 7 — Canal de Comunicação: Modelo Abstrato e Realidade Física}
% ============================================================

\subsection*{Canal ideal versus canal real}

No modelo abstrato, um canal de comunicação é um mapeamento matemático entre sinal
transmitido $x(t)$ e sinal recebido $y(t)$. O canal AWGN ideal satisfaz $y(t) = x(t) + n(t)$,
onde $n(t)$ é ruído gaussiano branco de densidade espectral $N_0/2$. Esse modelo é
tratável analiticamente e fornece limites teóricos (capacidade de Shannon), mas \textit{não
existe} fisicamente.

O canal real é um sistema dinâmico, não-linear, variante no tempo e sujeito a múltiplos
fenômenos simultâneos. A discrepância entre modelo e realidade é a principal fonte de
complexidade no projeto de sistemas de comunicação.

\subsection*{Impactos dos fenômenos físicos}

\textbf{Atenuação:} A amplitude do sinal decresce com a distância (lei inversa do quadrado
para propagação no espaço livre; lei exponencial para meios guiados). No receptor, isso
reduz a SNR efetiva. Sem controle de ganho automático (AGC), o receptor pode operar em
saturação ou em sensibilidade insuficiente dependendo da distância.

\textbf{Ruído:} O ruído térmico ($kTBF$, onde $F$ é a figura de ruído do receptor) é
inevitável. Fontes externas adicionais incluem ruído atmosférico, interferência de outros
sistemas e ruído de fase dos osciladores locais. O ruído impõe um limite fundamental à
distância máxima de operação sem amplificação.

\textbf{Interferência:} Diferentemente do ruído (estocástico), a interferência tem
estrutura determinística mas variável: outros transmissores no mesmo canal ou em canais
adjacentes. Em sistemas OFDM (Wi-Fi, LTE), a interferência inter-subportadora (ICI)
surge quando há erro de frequência entre transmissor e receptor.

\textbf{Atraso:} O atraso de propagação é fixo ($d/v_p$, onde $v_p$ é a velocidade de
fase no meio). Mais problemático é o atraso de grupo variável com a frequência (dispersão),
que provoca distorção de fase e, consequentemente, ISI.

\subsection*{Estratégias de mitigação pela engenharia}

A engenharia mitiga esses efeitos em múltiplos níveis: codificação de canal com FEC
(turbo codes, LDPC, polar codes) para ganho de codificação contra ruído; equalização
adaptativa (LMS, RLS) para compensação de distorção linear; técnicas de diversidade
(espacial, temporal, de frequência) para combater desvanecimento; OFDM com prefixo
cíclico para eliminar ISI em canais multipercurso; e amplificadores com controle automático
de ganho (AGC) para compensar atenuação variável.

% ============================================================
\section*{Questão 8 — Representação, Codificação e o Canal}
% ============================================================

\subsection*{Por que a representação deve ser compatível com o canal}

Toda transmissão de informação exige a escolha de uma \textbf{representação}: como os
símbolos do alfabeto da fonte são mapeados em sinais físicos. Essa escolha não é neutra ---
ela determina o espectro do sinal transmitido, a eficiência energética e a robustez a
degradações do canal.

A codificação de linha, por exemplo, define a densidade de transições no sinal transmitido.
A codificação NRZ-L tem espectro com componente DC significativa, tornando-a inadequada
para canais passa-banda ou com acoplamento capacitivo (que bloqueia DC). A codificação
Manchester elimina a componente DC ao custo de dobrar a largura de banda necessária.
A codificação 4B5B mapeia 4 bits em 5 bits com densidade de transições garantida,
permitindo recuperação de relógio sem ampliar tanto o espectro quanto Manchester.

\subsection*{Relação com sincronismo}

O sincronismo de bit (recuperação de relógio) depende da densidade de transições no sinal.
Um sinal NRZ com longas sequências de bits idênticos não apresenta transições, impossibilitando
que o receptor sincronize o relógio local ao fluxo de dados. Codificações que garantem
densidade mínima de transições (como 8B10B em Gigabit Ethernet) resolvem esse problema
de forma elegante, introduzindo apenas 25\% de overhead de banda.

\subsection*{Relação com energia transmitida}

A energia por bit $E_b$ é o produto da potência de transmissão pelo período de bit $T_b = 1/R_b$.
Aumentar a taxa $R_b$ sem aumentar a potência reduz $E_b$ proporcionalmente, degradando a
razão $E_b/N_0$ e, portanto, aumentando a BER. A escolha de modulação (BPSK, QPSK, QAM)
define o compromisso entre eficiência espectral e eficiência energética: esquemas de ordem
superior (256-QAM) transmitem mais bits por símbolo, mas exigem SNR muito maior para a
mesma BER.

\subsection*{Robustez à distorção}

Sinais com amplo espectro são mais vulneráveis à resposta em frequência não plana do canal
(distorção). Uma solução é o \textit{pulse shaping} com filtros de Nyquist (cosseno elevado),
que concentra a energia do pulso em uma banda controlada e minimiza ISI no receptor
amostrador. Outra solução é a multiplexação por subportadoras (OFDM), que divide o canal
seletivo em frequência em $N$ subcanais planos, cada um modulado independentemente, tornando
o sistema resiliente à dispersão sem necessidade de equalizadores complexos por subportadora.

% ============================================================
\section*{Questão 9 — Análise Crítica do Modelo OSI}
% ============================================================

\subsection*{Vantagens do modelo}

\textbf{Vantagem 1 --- Abstração e modularidade:} O modelo OSI provê um vocabulário comum
e limites de responsabilidade bem definidos, permitindo que equipes de hardware (camadas 1-2)
e de software (camadas 3-7) trabalhem independentemente. A substituição de Ethernet por
Wi-Fi na camada 2, por exemplo, não exige modificação dos protocolos IP ou TCP nas camadas
superiores, pois a interface entre camadas é estável. Essa propriedade é o fundamento da
evolução tecnológica das redes: sem a abstração de camadas, cada novo meio físico exigiria
redesign completo da pilha de software.

\textbf{Vantagem 2 --- Ferramenta de diagnóstico:} O modelo estrutura a análise de falhas
em redes: a metodologia de \textit{troubleshooting} em camadas (verificar física → enlace →
rede → transporte → aplicação) é universalmente ensinada e praticada. Ferramentas como
Wireshark organizam capturas de pacotes exatamente segundo as camadas OSI, tornando a
análise intuitiva.

\subsection*{Limitações do uso estritamente teórico}

\textbf{Limitação 1 --- Não reflete a arquitetura real:} A pilha TCP/IP dominante na
Internet não mapeia biunivocamente sobre OSI. As camadas de Sessão e Apresentação são
praticamente inexistentes como entidades separadas: o TLS opera entre Transporte e
Aplicação, o HTTP trata de sessão e representação de dados simultaneamente. Tentar
encaixar TCP/IP no OSI gera confusão pedagógica e análises incorretas.

\textbf{Limitação 2 --- Ignorar otimizações cross-layer:} Sistemas modernos de comunicação
sem fio (LTE, 5G, Wi-Fi 6) utilizam informações de camadas inferiores (qualidade do canal,
medidas de RSSI) para tomar decisões em camadas superiores (escalonamento de pacotes,
seleção de modulação). Esse \textit{cross-layer design} é deliberado e necessário para
eficiência, mas viola formalmente a separação de camadas do modelo. O OSI, aplicado
rigidamente, desencorajaria tais otimizações.

\subsection*{Conclusão: por que o modelo permanece relevante}

O modelo OSI permanece relevante não como especificação de implementação, mas como
\textbf{framework conceitual de decomposição de problemas}. Em ensino, fornece
estrutura para apresentar a complexidade das redes de forma progressiva. Em projeto,
orienta a análise de interoperabilidade e a detecção de responsabilidades mal alocadas.
A chave é utilizá-lo como ferramenta intelectual flexível, não como dogma arquitetural:
reconhecer onde o modelo se aplica bem (separação física/enlace/rede) e onde exige
adaptação (camadas superiores em sistemas modernos).

% ============================================================
\section*{Questão 10 — Texto Integrador: Uma Cadeia Lógica Única}
% ============================================================

A comunicação de dados, em sua essência, é o problema de transportar informação de uma fonte
para um destino de forma fiel, eficiente e confiável. Compreender por que esse problema é
difícil --- e como a engenharia o aborda --- exige a integração dos cinco pilares conceituais
estudados nos Capítulos 1 a 5: fundamentos do sistema de comunicação, arquitetura em camadas,
teoria de sinais, meios físicos e teoria do canal.

O ponto de partida é o modelo sistêmico (Cap.~1): a comunicação é uma cadeia de elementos
interdependentes. Esse modelo não é apenas descritivo --- é normativo. Ele estabelece que
o desempenho global é limitado pelo elo mais fraco e que a otimização local de um componente
sem considerar os demais pode ser ineficaz ou contraproducente. Um transmissor de altíssima
eficiência energética é inútil se o meio físico introduzir atenuação que o receptor não
consegue compensar.

O meio físico (Cap.~4) é o elemento onde as restrições mais fundamentais residem. Suas
propriedades --- atenuação, banda, susceptibilidade a EMI --- não são configuráveis por
software. São determinadas pela física dos materiais e pela geometria do enlace. A escolha
do meio (par trançado, coaxial, fibra óptica, canal de rádio) define os parâmetros dentro
dos quais todos os outros elementos devem operar.

Para que a informação da fonte chegue ao destino através desse meio, ela deve ser convertida
em um sinal compatível com as propriedades físicas do canal. É aqui que a teoria de sinais
(Cap.~3) se torna indispensável: a análise de Fourier revela que qualquer sinal com taxa de
variação finita ocupa uma largura de banda finita. A escolha de representação, modulação e
codificação de linha define não apenas o espectro ocupado, mas a sensibilidade do sinal à
distorção, ao ruído e às limitações do meio.

O canal real (Cap.~5) impõe sobre o sinal transmitido efeitos que o modelo simplificado não
prevê: atenuação variável, distorção, ruído e interferência. A teoria do canal quantifica
esses efeitos e estabelece limites fundamentais --- a capacidade de Shannon --- que nenhum
sistema pode violar. A engenharia de canal (FEC, equalização, diversidade) opera no espaço
entre o canal ideal e o canal real, tentando aproximar o desempenho prático ao limite teórico.

A arquitetura em camadas (Cap.~2) é a solução de engenharia para gerenciar a complexidade
resultante dessa cadeia. Cada camada resolve um subproblema específico e abstrai os demais:
a camada Física lida com sinais no meio; o Enlace provê confiabilidade local; a Rede trata
do roteamento global. Essa separação permite que avanços em um nível --- como o desenvolvimento
de fibras de menor atenuação --- beneficiem os protocolos de camadas superiores sem exigir
redesign da pilha completa.

A interdependência entre essas dimensões é a conclusão central: um problema de alta taxa de
erros em uma rede pode ter origem no meio físico (atenuação excessiva), no sinal (modulação
inadequada ao perfil de ruído), no canal (interferência não prevista no projeto), no enlace
(ausência de FEC) ou na rede (fragmentação excessiva causando RTT alto). Isolar o diagnóstico
a um único nível sem considerar a cadeia completa leva a soluções paliativas que frequentemente
deslocam o problema sem resolvê-lo. É por isso que a formação em comunicação de dados exige
o domínio integrado --- e não fragmentado --- desses cinco pilares.

% ============================================================
% Bibliografia baseada no Plano de Ensino 2026/1 e Programa 445AH-04
\begin{thebibliography}{9}

% ---- Bibliografia Principal ----
\bibitem{rochol2012}
J. Rochol, \textit{Comunicação de Dados}, vol.~22, Série Livros Didáticos Informática UFRGS.
Porto Alegre: Bookman, 2012.

% ---- Complementar utilizada (Q2: arquitetura de camadas e sistemas digitais) ----
\bibitem{sklar2001}
B. Sklar, \textit{Digital Communications: Fundamentals and Applications}, 2.~ed.
Upper Saddle River: Prentice Hall, 2001.

% ---- Complementar utilizada (Q3, Q5, Q7, Q8: modulação, ruído, canal, BER) ----
\bibitem{haykin2004}
S. Haykin, \textit{Sistemas de Comunicação: Analógicos e Digitais}, 4.~ed.
Porto Alegre: Bookman, 2004.

% ---- Complementar utilizada (Q3, Q5, Q8: Shannon, modulação digital, ISI, OFDM) ----
\bibitem{proakis2008}
J. G. Proakis, \textit{Digital Communications}, 5.~ed. New York: McGraw-Hill, 2008.

\end{thebibliography}

\end{document}